\section{Independencia de los axiomas de Huntington de 1904 (Opcional)}
Esta sección explora la independencia lógica de los postulados propuestos por Edward V. Huntington en 1904.\footnote{Edward V. Huntington, \textit{Sets of Independent Postulates for the Algebra of Logic}, Transactions of the American Mathematical Society, Vol. 5, No. 3 (Jul., 1904), pp. 288-309.}

\desplegable{independencia}{
\newcommand{\fromto}{\ensuremath{\longrightarrow}}
\newcommand{\bfs}[1]{\ensuremath{\boldsymbol{#1}}}
\newcommand{\bfsf}[1]{\ensuremath{\bfs{\mathsf{#1}}}}
\newcommand{\funcsum}{\ensuremath{\bfs{\mathsf{s}}}}
\newcommand{\funcsumat}[2]{\ensuremath{\funcsum\left({#1,#2}\right)}}
\newcommand{\funcprod}{\ensuremath{\bfs{\mathsf{p}}}}
\newcommand{\funcprodat}[2]{\ensuremath{\funcprod\left({#1,#2}\right)}}
\newcommand{\mathdef}{\ensuremath{\stackrel{\mathsf{def}}{{}={}}}}
\newcommand{\pfmathdef}[2]{\ensuremath{{#1}{\mathdef}{#2}}}
\newcommand{\setB}{\ensuremath{\mathbb{B}}}
\newcommand{\bfsf}[1]{\mathbf{\mathsf{#1}}}
\newcommand{\bfs}[1]{\mathbf{#1}}
\newcommand{\funccomp}{\ensuremath{\bfsf{c}}}
\newcommand{\fromto}{\ensuremath{\rightarrow}}
\newcommand{\setB}{\ensuremath{\mathbb{B}}}




	Para esto demostrar que los axiomas antes dados son independientes entre
sí, primero daremos un modelo consistente y sencillo, en el que haremos
variaciones y obtendremos modelos (ejemplos) de sistemas que cumplan todos
los axiomas excepto uno de ellos. Si se logra quedará claro que no podemos
deducir el axioma fallido del resto de axiomas que sí que se cumplen: el
axioma fallido es lógicamente independiente del resto de axiomas. Esto es
fácil de conseguir, comenzando con modelos de álgebras con base en un
conjunto de dos elementos \({\setB}_2 = \set{0,1}\). En esto copiamos los
modelos dados por Huntington en su artículo de 1904.



\subsection{Modelo mínimo de álgebra de Boole.}



	Definimos:
\begin{flalign}
{\setB}_2 &\mathdef \set{0,1} = \set{0} \cup \set{1} \\
0 &\mathdef \emptyset \\
1 &\mathdef \set{\emptyset} \cup \set{ \set{ \emptyset } } = \set{ \emptyset,
\set{ \emptyset } } \\
0 &\;\in\; 1 \\
0 &\;\subsetneq\; 1 \\
0 &\;\neq\; 1 \\
\end{flalign}
\begin{flalign}
\bfsf{s} &:  {\setB}_2 \times {\setB}_2{\qquad}\fromto{\qquad} {\setB}_2\\
\end{flalign}
\begin{flalign}
\bfsf{s} &:: \left({0,0}\right) {\qquad} \mapsto {\qquad} 0\\
\bfsf{s} &:: \left({0,1}\right) {\qquad} \mapsto {\qquad} 1\\
\bfsf{s} &:: \left({1,0}\right) {\qquad} \mapsto {\qquad} 1\\
\bfsf{s} &:: \left({1,1}\right) {\qquad} \mapsto {\qquad} 1
\end{flalign}
\begin{flalign}
\bfsf{p} &:  {\setB}_2 \times {\setB}_2{\qquad}\fromto{\qquad} {\setB}_2\\
\end{flalign}
\begin{flalign}
\bfsf{p} &:: \left({0,0}\right) {\qquad} \mapsto {\qquad} 0\\
\bfsf{p} &:: \left({0,1}\right) {\qquad} \mapsto {\qquad} 0\\
\bfsf{p} &:: \left({1,0}\right) {\qquad} \mapsto {\qquad} 0\\
\bfsf{p} &:: \left({1,1}\right) {\qquad} \mapsto {\qquad} 1
\end{flalign}
\begin{flalign}
\bfsf{c} &:  {\setB}_2 {\qquad} \fromto {\qquad} {\setB}_2
\end{flalign}
\begin{flalign}
\bfsf{c} &:: 0 {\qquad} \mapsto {\qquad} 1\\
\bfsf{c} &:: 1 {\qquad} \mapsto {\qquad} 0
\end{flalign}



	El conjunto \({\setB}_2\) es conjunto en \textbf{ZFS}, desde el momento
que la el axioma de unión nos asegura que \({\setB}_2\) es conjunto unión de
dos conjuntos de un solo elemento, \({\setB}_2 \mathdef { \set{0} } \cup {
\set{1} } \). Los elementos \(0 \mathdef \emptyset\) y  \(1 \mathdef \set{
\emptyset, \set{ \emptyset } } \). De nuevo para definir \(1\) como conjunto
necesitamos el axioma de unión (o de pares no ordenados) de \textbf{ZFS},
siendo \(1 \mathdef \set{0}\cup\set{\set{0}}\). Podemos ver que \(0 \cap 1 =
\emptyset \) por lo que \(0 \neq 1\). También \(0 \in 1\). Todo este párrafo
constituye la satisfacción de los requerimientos \ref{axm:H0}.



	Tal como hemos definido nuestro modelo de álgebra de Boole, lo primero
que queda claro es que \(\funcsum\) y \(\funcprod\) son funciones binarias
bien definidas, y esta última es además biyectiva. Luego los axiomas
\ref{axm:H1s} y \ref{axm:H1p} quedan satisfechos. 



	Los axiomas \ref{axm:H2s} y \ref{axm:H2p} (existencia del elemento
neutro) quedan directamente satisfechos por simple inspección de las tablas. 



	Para el axioma \ref{axm:H2s} observamos que \(\funcsumat{0}{0} = 0\) y
\(\funcsumat{0}{1} = \funcsumat{1}{0} = 1\) nos muestra el elemento neutro de
la suma, el elemento \(0\).



	Para el axioma \ref{axm:H2p} observamos que \( \funcprodat{1}{0} =
\funcprodat{0}{1} = 0\) y \(\funcprodat{1}{1} = 1\) nos muestra el elemento
neutro del producto, el elemento \(1\).



	Para los axiomas de conmutatividad solo hay que observar en la definición
de las funciones binarias \(\funcsum,\funcprod : \setB_2 \times \setB_2
\fromto \setB_2\), que \(\funcsumat{0}{1} = \funcsumat{1}{0} = 1\) y se
cumple \ref{axm:H3s}, que \(\funcprodat{0}{1} = \funcprodat{1}{0} = 0\) y se
cumple \ref{axm:H3p}. 



	Para las distribuciones de una operación interna sobre otra elegiré
construir unas tablas que muestren la igualdades necesarias.



	Primero veremos la distribución de la suma sobre el producto, o, lo que
es lo mismo, su inversa, sacar factor común. 



	Lo haremos rellenando y comparando la siguiente tabla de verdad.



	El método de relleno es fácil. Primero rellenamos las tres primeras
columnas, corresponden a las variables independientes, que al ser tres con
dos valores posibles en cada ocasión, son un total de ocho combinaciones
distintas. Esas columnas están marcadas con \(1\) en la última fila (fila que
se ve separada). A continuación rellenamos las columnas correspondientes a
una operación binaria interna de dos variables independientes ya rellenas.
Para rellenarlas, solo tenemos que buscar los valores en las tablas dadas
para las operaciones internas en ka definición de este álgebra de Boole.
Estas están marcadas con un \(2\) en la última línea, la línea separada.
Corresponden con la cuarta columna, la suma \( \funcsumat{ \textsf{b} } {
\textsf{c} } \), y con las columnas siete y ocho: la ocho es un producto
\(\funcprodat{\textsf{a}}{\textsf{c}}\) y la siete otro \( \funcprodat{
\textsf{a} } { \textsf{b} } \). Las columnas centrales, la cinco y la seis
corresponden a todas las valoraciones posibles (y en el mismo orden de valor
de las variables independientes) de las dos expresiones que queremos
comparar, las afirmadas por el postulado \ref{axm:H4s}. Podemos ver que ambas
columnas son idénticas, luego se cumple el citado postulado.



\begin{flalign}
\begin{matrix}
\hline
\textsf{a} & \textsf{b} & \textsf{c} &
\funcsumat{\textsf{b}}{\textsf{c}} &
\funcprodat{\textsf{a}}{\funcsumat{\textsf{b}}{\textsf{c}}} &
\funcsumat{\funcprodat{\textsf{a}} {\textsf{b}}} {\funcprodat{\textsf{a}}
{\textsf{c}}}
&
{\funcprodat{\textsf{a}}{\textsf{b}}}
&
{\funcprodat{\textsf{a}}{\textsf{c}}}
\\
\hline
0 & 0 & 0 & 0 & 0 & 0 & 0 & 0 \\
0 & 0 & 1 & 1 & 0 & 0 & 0 & 0 \\
0 & 1 & 0 & 1 & 0 & 0 & 0 & 0 \\
0 & 1 & 1 & 1 & 0 & 0 & 0 & 0 \\
1 & 0 & 0 & 0 & 0 & 0 & 0 & 0 \\
1 & 0 & 1 & 1 & 1 & 1 & 0 & 1 \\
1 & 1 & 0 & 1 & 1 & 1 & 1 & 0 \\
1 & 1 & 1 & 1 & 1 & 1 & 1 & 1 \\
\hline
1 & 1 & 1 & 2 & 3 & 3 & 2 & 2 \\ 
\hline
\end{matrix}
\end{flalign}



	Por último veremos la distribución del producto sobre la suma, o, lo que
es lo mismo, su inversa, sacar sumando común. El método es en todo idéntico
al seguido para la otra distributiva.



\begin{flalign}
\begin{matrix}
\hline
\textsf{a} & \textsf{b} & \textsf{c} &
\funcprodat{\textsf{b}}{\textsf{c}} &
\funcsumat {\textsf{a}}
           {\funcprodat{\textsf{b}}{\textsf{c}}} &
\funcprodat{\funcsumat{\textsf{a}}{\textsf{b}}}
           {\funcsumat{\textsf{a}}{\textsf{c}}}
&
{\funcsumat{\textsf{a}}{\textsf{b}}}
&
{\funcsumat{\textsf{a}}{\textsf{c}}}
\\
\hline
0 & 0 & 0 & 0 & 0 & 0 & 0 & 0 \\
0 & 0 & 1 & 0 & 0 & 0 & 0 & 1 \\
0 & 1 & 0 & 0 & 1 & 0 & 1 & 0 \\
0 & 1 & 1 & 1 & 1 & 1 & 1 & 1 \\
1 & 0 & 0 & 0 & 1 & 1 & 1 & 1 \\
1 & 0 & 1 & 0 & 1 & 1 & 1 & 1 \\
1 & 1 & 0 & 0 & 1 & 1 & 1 & 1 \\
1 & 1 & 1 & 1 & 1 & 1 & 1 & 1 \\
\hline
1 & 1 & 1 & 2 & 3 & 3 & 2 & 2 \\ 
\hline
\end{matrix}
\end{flalign}




	Para el axioma \ref{axm:H5} observamos que \(\funccomplat{0} = 1\) y
\(\funccomplat{1} = 0\) y por inspección en las tablas vemos que
\(\funcsumat{0}{\funccomplat{0}} = \funcsumat{0}{1} = 1\) y que
\(\funcprodat{0}{\funccomplat{0}} = \funcprodat{0}{1} = 0\) cumpliendo para
el elemento \(0\) se cumple que \(\exists 1 = \funccomplat{0} \in \setB_2 \)
\(\funcsumat{1}{\funccomplat{1}} = \funcsumat{1}{0} = 1\), que coincide con
la ecuación segunda (ecuación: 2.15) de \ref{eqn:H5} y
\(\funcprodat{1}{\funccomplat{1}}=\funcprodat{1}{0} = 0\) que coincide con la
tercera ecuación (ecuación: 2.16) de \ref{eqn:H5}, e igualmente para el
elemento \( 1 \), \(\exists 0 = \funccomplat{1} \in \setB_2 \) \(
\funcsumat{0}{\funccomplat{0}} = \funcsumat{0}{1} = 1\) (ecuación: 2.15) y
\(\funcprodat{0}{\funccomplat{0}} =\funcprodat{0}{1} = 0\) (ecuación: 2.16)
de \ref{eqn:H5}.




	Como las ecuaciones 2.15 y 2.16 se cumplen para \(0\) y \(1\), esto es,
\(\forall x \in \setB_2\) como se requiere en la ecuación 2.14, queda
satisfecho el postulado \ref{axm:H5}.



\subsection{Independencia de la suma está siempre definida.}



Modelo en el que solo falla \ref{axm:H1s}, esto es, que la operación suma
no es operación interna: no es función, pero en el sentido que \(1 + 1 \notin
\setB\).


\[
    \begin{matrix}
        {\cdot + \cdot} & \big{\vert} & {0} & {1} & \big{\vert} \\
        \hline
        {0}             & \big{\vert} & {0} & {1} & \big{\vert} \\
        {1}             & \big{\vert} & {1} & {x} & \big{\vert} \\
        \hline
    \end{matrix}
    \qquad\qquad
    \begin{matrix}
        {\cdot * \cdot} & \big{\vert} & {0} & {1} & \big{\vert} \\
        \hline
        {0}                   & \big{\vert} & {0} & {0} & \big{\vert} \\
        {1}                   & \big{\vert} & {0} & {1} & \big{\vert} \\
        \hline
    \end{matrix}
\]
%% (H1s) 
%% 0 + 0 = 0; 0 + 1 = 1; 1 + 0 = 1; 1 + 1 = x; 
%% 0 * 0 = 0; 0 * 1 = 0; 1 * 0 = 0; 1 * 1 = 1; 



	Al ser \(1 + 1 = x\), dónde \(x\) es ningún elemento de $\setB$, o dicho
de otro modo, \(1 + 1 \notin \setB\), vemos que sigue existiendo el elemento
neutro de la suma, \(0\), que la suma sigue siendo conmutativa y las
distributivas siguen valiendo mientras la suma tenga sentido (mientras no
aparezca \(x\)). El complementario de \(0\) es \(1\) y el de \(1\) es \(0\).



\subsection{Independencia de la unicidad de la definición de la suma.}

Modelo en el que no se cumple \ref{axm:H1s}, en el sentido que \(1 + 1
\mapsto 1\) y \(1 + 1 \mapsto 0\).


\[
    \begin{matrix}
        {\cdot + \cdot} & \big{\vert} & {0} & {1} & \big{\vert} \\
        \hline
        {0}             & \big{\vert} & {0} & {1} & \big{\vert} \\
        {1}             & \big{\vert} & {1} & {x=0 \vee x=1} & \big{\vert} \\
        \hline
    \end{matrix}
    \qquad\qquad
    \begin{matrix}
        {\cdot * \cdot} & \big{\vert} & {0} & {1} & \big{\vert} \\
        \hline
        {0}                   & \big{\vert} & {0} & {0} & \big{\vert} \\
        {1}                   & \big{\vert} & {0} & {1} & \big{\vert} \\
        \hline
    \end{matrix}
\]
%% (H1s) 
%% 0 + 0 = 0; 0 + 1 = 1; 1 + 0 = 1; 1 + 1 \mapsto 0,1; 
%% 0 * 0 = 0; 0 * 1 = 0; 1 * 0 = 0; 1 * 1 = 1; 



Exactamente igual que en la versión anterior, se cumplen todos los demás
postulados, sea cual sea el valor tomado para \(1 + 1\).



\subsection{Independencia de el producto está siempre definido.}



	Modelo en el que solo falla \ref{axm:H1p}, esto es, que la operación
producto no es operación interna: no es función. Este es el caso en que
\( 0 \cdot 0 \notin \setB \). Ponemos en ese caso \(0 \cdot 0 = x\) pero
igualmente podríamos haber dejado en blanco ese lugar.



\[
    \begin{matrix}
        {\cdot + \cdot} & \big{\vert} & {0} & {1} & \big{\vert} \\
        \hline
        {0}             & \big{\vert} & {0} & {1} & \big{\vert} \\
        {1}             & \big{\vert} & {1} & {1} & \big{\vert} \\
        \hline
    \end{matrix}
    \quad
    \begin{matrix}
        {\cdot * \cdot} & \big{\vert} & {0} & {1} & \big{\vert} \\
        \hline
        {0}                   & \big{\vert} & {x} & {0} & \big{\vert} \\
        {1}                   & \big{\vert} & {0} & {1} & \big{\vert} \\
        \hline
    \end{matrix}
\]
%% (H1p) 
%% 0 + 0 = 0; 0 + 1 = 1; 1 + 0 = 1; 1 + 1 = 1; 
%% 0 * 0 = x; 0 * 1 = 0; 1 * 0 = 0; 1 * 1 = 1; 



	Como es fácil ver, existe elemento neutro tanto de la suma como del
producto. Ambas operaciones son conmutativas y se dan las dos propiedades de
distribución, siempre que tenga sentido el producto, esto es, no aparezca \(0 \cdot 0\). El complementario del \(0\) es \(\overline 0 = 1\) y el del
\(1\) es \(\overline 1 = 0\).



\subsection{Independencia de la unicidad de la definición del producto.}



Modelo en el que solo falla \ref{axm:H1s}, esto es, en el sentido que \(0
\cdot 0 \mapsto 0\) y también \(0 \cdot 0 \mapsto 1\).



\[
    \begin{matrix}
        {\cdot + \cdot} & \big{\vert} & {0} & {1} & \big{\vert} \\
        \hline
        {0}             & \big{\vert} & {0} & {1} & \big{\vert} \\
        {1}             & \big{\vert} & {1} & {1} & \big{\vert} \\
        \hline
    \end{matrix}
    \quad
    \begin{matrix}
        {\cdot * \cdot} & \big{\vert} & {0} & {1} & \big{\vert} \\
        \hline
        {0}                   & \big{\vert} & {x=0 \vee x=1} & {0} & \big{\vert} \\
        {1}                   & \big{\vert} & {0} & {1} & \big{\vert} \\
        \hline
    \end{matrix}
\]
%% (H2s) 
%% 0 + 0 = 0  ; 0 + 1 = 1; 1 + 0 = 1; 1 + 1 = 1; 
%% 0 * 0 = 0,1; 0 * 1 = 0; 1 * 0 = 0; 1 * 1 = 1;



\subsection{Independencia de la existencia de elemento neutro de la suma.}



Modelo en el que solo falla \ref{axm:H2s}, esto es, la existencia de
elemento neutro en la operación binaria interna suma.



\[
    \begin{matrix}
        {\cdot + \cdot} & \big{\vert} & {0} & {1} & \big{\vert} \\
        \hline
        {0}             & \big{\vert} & {0} & {0} & \big{\vert} \\
        {1}             & \big{\vert} & {0} & {0} & \big{\vert} \\
        \hline
    \end{matrix}
    \quad
    \begin{matrix}
        {\cdot * \cdot} & \big{\vert} & {0} & {1} & \big{\vert} \\
        \hline
        {0}                   & \big{\vert} & {0} & {0} & \big{\vert} \\
        {1}                   & \big{\vert} & {0} & {1} & \big{\vert} \\
        \hline
    \end{matrix}
\]
%% II.a H2s 
%% 0 + 0 = 0; 0 + 1 = 0; 1 + 0 = 0; 1 + 1 = 0; 
%% 0 * 0 = 0; 0 * 1 = 0; 1 * 0 = 0; 1 * 1 = 1;


	Primero, está claro que las operaciones internas están bien construidas
en cuanto a que son funciones.


	Queda claro que no hay elemento neutro de la suma pues \(\forall xy \in
\setB \quad x + y = 0\).


	Segundo, existe el elemento neutro del producto: \(0 \cdot 1 = 1 \cdot 0
= 0 \) y \( 1 \cdot 1 = 1 \). Luego se cumple \ref{aciom:H2p}.



	La conmutatividad se hace patente al ver las diagonales inversas de las
tablas de operación, que muestran un único valor. \( 0 + 1 = 1 + 0 = 0\) y
\(0 \cdot 1 = 1 \cdot 0 = 0\). Se cumplen \ref{axm:H3s} y \ref{axm:H3p}. 



	En cuanto a la distribución del producto sobre la suma, veamos si podemos
comprobarla de forma sencilla: \( a \cdot ( b + c ) = (a \cdot b) + (a \cdot
c) \). Sabemos que \(b + c = 0\) siempre, y que, pongamos que \(a \cdot b = x
\in \setB\) y que \(a \cdot c = y \in \setB\). Ahora bien \(x + y = 0\). Así
que todo lo que tenemos que probar es que \(a \cdot 0 = 0\), pero esto es
claro en la table del producto. Luego se cumple \ref{axm:H4s}.



	Ahora la distribución de la suma sobre el producto. \( a + ( b \cdot c )
= (a + b) \cdot (a + c) \). Sabemos que \(a + (b \cdot c) = 0\) siempre, y
que, pongamos que \( a + b = 0 \) y que \(a + c = 0 \). Ahora bien \(0 \cdot
0 = 0\). Luego se cumple \ref{axm:H4p}.


	Nos queda encontrar un complemento para el \(0\). Pero encontrar el
complemento solo tiene sentido si existen los dos elementos neutros, pero en
este caso no existe el neutro de la suma: no tiene sentido buscar el
complementario.



\subsection{Independencia de la existencia de elemento neutro del producto.}



	Exponemos un modelo en el que solo falla \ref{axm:H2p}, esto es, la existencia de elemento neutro en la operación binaria interna producto.



\[
    \begin{matrix}
        {\cdot + \cdot} & \big{\vert} & {0} & {1} & \big{\vert} \\
        \hline
        {0}             & \big{\vert} & {0} & {1} & \big{\vert} \\
        {1}             & \big{\vert} & {1} & {1} & \big{\vert} \\
        \hline
    \end{matrix}
    \quad
    \begin{matrix}
        {\cdot * \cdot} & \big{\vert} & {0} & {1} & \big{\vert} \\
        \hline
        {0}                   & \big{\vert} & {1} & {1} & \big{\vert} \\
        {1}                   & \big{\vert} & {1} & {1} & \big{\vert} \\
        \hline
    \end{matrix}
\]
%% (H2p) 
%% 0 + 0 = 0; 0 + 1 = 1; 1 + 0 = 1; 1 + 1 = 1; 
%% 0 * 0 = 1; 0 * 1 = 1; 1 * 0 = 1; 1 * 1 = 1;


	Este modelo es completamente simétrico al anterior, por lo que no nos
hará falta demostrar que se dan el resto de postulados. Solo hay que seguir
el esquema de la sub-sección anterior.



\subsection{Independencia de la conmutatividad de la suma.}



	Modelo en que la conmutatividad de la suma \ref{axm:H3s} no se dá, pero
si que se dan el resto de postulados.



\begin{equation}
    \begin{matrix}
        {\cdot + \cdot} & \big{\vert} & {0} & {1} & \big{\vert} \\
        \hline
        {0}             & \big{\vert} & {0} & {0} & \big{\vert} \\
        {1}             & \big{\vert} & {1} & {1} & \big{\vert} \\
        \hline
    \end{matrix}
    \quad
    \begin{matrix}
        {\cdot * \cdot} & \big{\vert} & {0} & {1} & \big{\vert} \\
        \hline
        {0}                   & \big{\vert} & {0} & {0} & \big{\vert} \\
        {1}                   & \big{\vert} & {0} & {1} & \big{\vert} \\
        \hline
    \end{matrix}
\end{equation}
%% (H3s) 
%% 0 + 0 = 0; 0 + 1 = 0; 1 + 0 = 1; 1 + 1 = 1; 
%% 0 * 0 = 0; 0 * 1 = 0; 1 * 0 = 0; 1 * 1 = 1;



	Que la suma y el producto están bien definidos como función es inmediato.



	Podemos comprobar que tanto \(\forall x \; x + 0 = x\) como que \(\forall
x \; x + 1 = x\). Esto quiere decir que la suma tiene dos neutros, el \(0\) y
el \(1\). Esto tendrá repercusiones sobre el elemento complementario. Podemos
ver por otra parte que \(1 + 0 \neq 0 + 1\), por lo que efectivamente la suma
no es conmutativa. Para el producto todo es idéntico, tiene elemento neutro
general y único \(0\) y \(1\cdot 0 = 0\cdot 1\) por lo que el producto es
conmutativo.



	Por otro lado, vemos que el complementario de \(0\) es \(1\), \(0 + 1 =
1\) -es neutro del producto- y \(0 * 1 = 0\) -es neutro de la suma-. Para ver
que \(1\) tiene complementario, hay que ir con más cuidado, \(1 + 1 = 1\) -
\(1\) es neutro del producto - y \(1 * 1 = 1\) - \(1\) es también neutro del
producto -. Luego existe complementario de \(1\) y es \(1\).



	Vamos ahora con las distribuciones:



	Nos preguntamos si se cumple \( x \cdot ( y + z ) = ( x \cdot y ) + ( x
\cdot z ) \) en el actual modelo. Si \(x = 0\) entonces \( 0 \cdot (x + y) =
0 \) y la parte derecha de la igualdad se resuelve en \( ( 0 \cdot y ) + ( 0
\cdot z ) \) pero como \( 0 \cdot x = 0 \) obtenemos que la distribución
\ref{axm:H4s} será verdad para \(x = 0\) si \( 0 = 0 + 0 \), cosa que es
cierta. Si \( x = 1 \), la igualdad a verificar quedaría \( 1 \cdot ( y + z )
= ( 1 \cdot y ) + ( 1 \cdot z ) \) y por \ref{axm:H2p} queda \( y + z  =  y
+ z \) que no es más que la identidad lógica de la igualdad. Se satisface
\ref{axm:H4s}.



	Nos preguntamos por la satisfacción de \ref{axm:H4p} \( x + ( y \cdot z
) = ( x + y ) \cdot ( x + z ) \) en el actual modelo. Procedemos como en el
párrafo anterior, por casos. Si \(x = 0\) entonces \( x + ( y \cdot z ) = ( x
+ y ) \cdot ( x + z ) \) \(\Longrightarrow\) \(0 + ( y \cdot z ) = ( 0 + y )
\cdot ( 0 + z )\) \(\Longrightarrow\) \( 0 = 0 \cdot 0\) lo que es cierto.
Para el caso \(x = 1\),  obtenemos \( x + ( y \cdot z ) = ( x + y ) \cdot ( x
+ z ) \) \(\Longrightarrow\) \( 1 + ( y \cdot z ) = ( 1 + y ) \cdot ( 1 + z
)\) \(\Longrightarrow\) \( 1 = 1 \cdot 1 \). Por lo tanto también se verifica
\ref{axm:H4p}.



	Queda comprobado que la conmutatividad de la suma \ref{axm:H3s} es independiente del resto de postulados.



\subsection{Independencia de la conmutatividad del producto.}



	Modelo en que la conmutatividad del producto \ref{axm:H3p} no se da,
pero, si se satisfacen el resto de postulados.



\begin{equation}
    \begin{matrix}
        {\cdot + \cdot} & \big{\vert} & {0} & {1} & \big{\vert} \\
        \hline
        {0}             & \big{\vert} & {0} & {1} & \big{\vert} \\
        {1}             & \big{\vert} & {1} & {1} & \big{\vert} \\
        \hline
    \end{matrix}
    \quad
    \begin{matrix}
        {\cdot * \cdot} & \big{\vert} & {0} & {1} & \big{\vert} \\
        \hline
        {0}                   & \big{\vert} & {0} & {0} & \big{\vert} \\
        {1}                   & \big{\vert} & {1} & {1} & \big{\vert} \\
        \hline
    \end{matrix}
\end{equation}
%% H3p 
%% 0 + 0 = 0; 0 + 1 = 1; 1 + 0 = 1; 1 + 1 = 1; 
%% 0 * 0 = 0; 0 * 1 = 0; 1 * 0 = 1; 1 * 1 = 1;



	El modelo es totalmente simétrico al anterior por lo que se puede
comprobar todo siguiendo los mismos pasos que en la sub-sección anterior.



\subsection{Independencia de la distribución de la suma sobre el producto.}



	El postulado del título no se cumple. Existen valores del modelo \(
\exists x \in \setB \) que no cumplen \ref{axm:H4p}, \( x + ( y \cdot z )
\neq ( x + y ) \cdot (x + z) \).



\begin{equation}
    \begin{matrix}
        {\cdot + \cdot} & \big{\vert} & {0} & {1} & \big{\vert} \\
        \hline
        {0}             & \big{\vert} & {0} & {1} & \big{\vert} \\
        {1}             & \big{\vert} & {1} & {0} & \big{\vert} \\
        \hline
    \end{matrix}
    \quad
    \begin{matrix}
        {\cdot * \cdot} & \big{\vert} & {0} & {1} & \big{\vert} \\
        \hline
        {0}                   & \big{\vert} & {0} & {0} & \big{\vert} \\
        {1}                   & \big{\vert} & {0} & {1} & \big{\vert} \\
        \hline
    \end{matrix}
\end{equation}
%% (H4s) 
%% 0 + 0 = 0; 0 + 1 = 1; 1 + 0 = 1; 1 + 1 = 0;
%% 0 * 0 = 0; 0 * 1 = 0; 1 * 0 = 0; 1 * 1 = 1;



	Este modelo es virtualmente idéntico al cuerpo sobre \(\setZ\) de restos
módulo 2, \(\setZ/{\mod2}\). Ese es el cambio que se da en la suma. Ahora \(1
+ 1 = 0\).



	Se dan por lo tanto todos los teoremas conocidos incluida la distribución
\ref{axm:H4s} e incluso sabemos que las operaciones son asociativas.



	Pero en aritmética no se da la dualidad que buscamos y la distribución
del producto sobre la suma no es una verdad.



	Solo nos queda ver si cumple la existencia de elemento complementario que
se encuentra visualmente por inspección de las tablas: \(\overline{1} = 0\) y
\(\overline{0} = 1\).


	Nos queda verificar que efectivamente no se cumple la distribución de la suma sobre el producto (que intuitivamente vemos). Encontramos un caso, \( 1 + ( y \cdot z ) \neq ( 1 + y ) \cdot ( 1 + z ) \). Lo vemos a continuación.


\begin{flalign}
  x = 1 &{} \\
  1 + ( y \cdot z ) &= ( 1 + y ) \cdot ( 1 + z ) \\
  \overline{y \cdot z} &= \overline{y} \cdot \overline{z} \\
  y = 0 &{} \\
  \overline{0 \cdot z} &= \overline{0} \cdot \overline{z} \\
  \overline{0} &= 1 \cdot \overline{z} \\
  1 &= \overline{z} \\
  z &= 0\\
  x=1 \wedge y=0 \wedge z=1 &\Longrightarrow \neg\ref{axm:H4p}
\end{flalign}



	De dónde efectivamente este modelo no distribuye la suma sobre un
producto. Y la independencia de \ref{axm:H4p} queda probada.



\subsection{Independencia de la distribución del producto sobre la suma.}


\begin{equation}
    \begin{matrix}
        {\cdot + \cdot} & \big{\vert} & {0} & {1} & \big{\vert} \\
        \hline
        {0}             & \big{\vert} & {0} & {1} & \big{\vert} \\
        {1}             & \big{\vert} & {1} & {1} & \big{\vert} \\
        \hline
    \end{matrix}
    \quad
    \begin{matrix}
        {\cdot * \cdot} & \big{\vert} & {0} & {1} & \big{\vert} \\
        \hline
        {0}                   & \big{\vert} & {1} & {0} & \big{\vert} \\
        {1}                   & \big{\vert} & {0} & {1} & \big{\vert} \\
        \hline
    \end{matrix}
\end{equation}
%% (H4s) 
%% 0 + 0 = 0; 0 + 1 = 0; 1 + 0 = 0; 1 + 1 = 1; 
%% 0 * 0 = 1; 0 * 1 = 0; 1 * 0 = 0; 1 * 1 = 1;



	Completamente simétrico al caso anterior.



\subsection{Independencia de la existencia del elemento complementario.}



	Por último presentamos un modelo que no cumple la existencia del
complementario y si satisface el resto de postulados.



\begin{equation}
    \begin{matrix}
        {\cdot + \cdot} & \big{\vert} & {0} & {1} & \big{\vert} \\
        \hline
        {0}             & \big{\vert} & {0} & {1} & \big{\vert} \\
        {1}             & \big{\vert} & {1} & {1} & \big{\vert} \\
        \hline
    \end{matrix}
    \quad
    \begin{matrix}
        {\cdot * \cdot} & \big{\vert} & {0} & {1} & \big{\vert} \\
        \hline
        {0}                   & \big{\vert} & {0} & {1} & \big{\vert} \\
        {1}                   & \big{\vert} & {1} & {1} & \big{\vert} \\
        \hline
    \end{matrix}
\end{equation}
%% (H5) 
%% 0 + 0 = 0; 0 + 1 = 1; 1 + 0 = 1; 1 + 1 = 1;
%% 0 * 0 = 0; 0 * 1 = 1; 1 * 0 = 1; 1 * 1 = 1;



	Como podemos ver en este modelo la suma y el producto son idénticos, por
lo que el \(0\) es el único neutro, tanto de la suma como del producto. La
conmutatividad es evidente a primera vista. 


	Veamos la distribución del producto respecto de la suma. Com la suma y el
producto son idénticos, bastará con probar este caso únicamente.



\begin{equation}
x \cdot (y + z) = x + ( y + z ) = ( x + x ) + ( y + z ) = ( x + y ) + ( x + z
) = ( x + y ) \cdot ( x + z )
\end{equation}



	Solo queda probar que \(x = x + x\) (que se hace por simple inspección
visual de las tablas) y la propiedad asociativa para este modelo, que
tendremos que probar. Lo haremos mediante una tabla.



\begin{equation}
  \begin{matrix}
  \hline
  \mathsf{a} & \mathsf{b} & \mathsf{c} & \mathsf{b}\cdot\mathsf{c} &
  \mathsf{a}\cdot\left( \mathsf{b} \cdot \mathsf{c} \right) &
  \left(\mathsf{a}\cdot \mathsf{b}\right) \cdot \mathsf{c} &
  \mathsf{a}\cdot\mathsf{b}\\
  \hline
%%a   b   c   bc a(bc)(ab)c ab
  0 & 0 & 0 & 0 & 0    & 0 & 0 \\
  0 & 0 & 1 & 1 & 1    & 1 & 0 \\
  0 & 1 & 0 & 1 & 1    & 1 & 1 \\
  0 & 1 & 1 & 1 & 1    & 1 & 1 \\
  1 & 0 & 0 & 0 & 1    & 1 & 1 \\
  1 & 0 & 1 & 1 & 1    & 1 & 1 \\
  1 & 1 & 0 & 1 & 1    & 1 & 1 \\
  1 & 1 & 1 & 1 & 1    & 1 & 1 \\
  \hline
  1 & 1 & 1 & 2 & 3    & 3 & 2 \\
  \hline
  \end{matrix}
\end{equation}


}
